\scanpage{161}
% archive foliation: 160
section des idéaux premiers $\mathfrak j_y$ pour $y\in Y$. Il est clair que la relation $Y\subset Y'$ entraîne $\mathfrak j(Y)\supset\mathfrak j(Y')$ et que l'on a
\[
\mathfrak j(\bigcup_\lambda Y_\lambda)=\bigcap_\lambda\mathfrak j(Y_\lambda)\tag{1.1.3.1}
\]
pour toute famille $(Y_\lambda)$ de parties de X.
\inserted{Enfin $\boxed{\mathfrak j(\{x\})=\mathfrak j_x}$}
\margin{(sert pour 1.1.7. !)}

\textit{Proposition} \textbf{(1.1.4).} --- (i) \textit{Pour toute partie E de A, on a $\mathfrak j(V(E))=\mathrm r(E)$.}

(ii) \textit{Pour toute partie Y de X, $V(\mathfrak j(Y))=\overline Y$, adhérence de Y dans X.}

(i) est conséquence immédiate des définitions et de (1.1.1.1); d'autre part, $V(\mathfrak j(Y))$ est fermé et contient Y; inversement, si $Y\subset V(E)$, on a $f(y)=0$ pour tout $f\in E$ et tout $y\in Y$, donc $E\subset\mathfrak j(Y)$, $V(E)\supset V(\mathfrak j(Y))$, ce qui prouve (ii).

\textit{Corollaire} \textbf{(1.1.5).} --- \textit{Les parties fermées de $X=\operatorname{Spec}(A)$ et les idéaux de A égaux à leurs racines (autrement dit les intersections d'idéaux premiers) se correspondent biunivoquement par les applications décroissantes $Y\to\mathfrak j(Y)$, $\mathfrak a\to V(\mathfrak a)$; à la réunion $Y_1\cup Y_2$ de deux parties fermées correspond $\mathfrak j(Y_1)\cap\mathfrak j(Y_2)$, et à l'intersection d'une famille quelconque $(Y_\lambda)$ de parties fermées correspond la racine de la somme des $\mathfrak j(Y_\lambda)$.}

\textit{Corollaire} \textbf{(1.1.6).} --- \textit{Si A est un anneau noethérien, $X=\operatorname{Spec}(A)$ est un espace noethérien..}

On notera que la réciproque de ce corollaire est inexacte, comme le montre l'exemple d'un anneau intègre non noethérien ayant un seul idéal premier $\ne\{0\}$, par exemple un anneau de valuation non discrète de rang 1.

Comme exemple d'anneau A dont le spectre n'est pas un espace noethérien, on peut signaler l'anneau $\mathcal C(Y)$ des fonctions continues numériques sur un espace compact infini Y; on sait qu'en tant qu'ensemble, Y s'identifie à l'ensemble des idéaux maximaux de A, et il est facile de voir que la topologie induite sur Y par celle de $X=\operatorname{Spec}(A)$ est la topologie initiale de Y. Comme Y n'est pas un espace noethérien, il en est donc de même de X.

\textit{Corollaire} \textbf{(1.1.7).} --- \textit{Pour tout $x\in X$, l'adhérence de $\{x\}$ est l'ensemble des $y\in X$ tels que $\mathfrak j_x\subset\mathfrak j_y$. Pour que $\{x\}$ soit fermé, il faut et il suffit que $\mathfrak j_x$ soit maximal.}

\textit{Corollaire} \textbf{(1.1.8).} --- \textit{L'espace $X=\operatorname{Spec}(A)$ est un espace de Kolmogoroff.}

En effet, si $x,y$ sont deux points distincts de X, on a, soit $\mathfrak j_x\not\subset\mathfrak j_y$, soit $\mathfrak j_y\not\subset\mathfrak j_x$, donc l'un des points $x,y$ n'appartient pas à l'adhérence de l'autre.

\textbf{(1.1.9)} D'après la prop. (1.1.\deleted{1}\inserted{2}, (iv)), pour deux éléments $f,g$ de A, on a
\[
D(fg)=D(f)\cap D(g)\tag{1.1.9.1}
\]
Notons aussi que la relation $D(f)=D(g)$ signifie, d'après la prop. (1.1.4, (i)) et la prop. (1.1.1, (v)) que $\mathrm r(f)=\mathrm r(g)$, ou encore que les idéaux premiers minimaux contenant $(f)$ et $(g)$ sont les mêmes; en particulier, il en est ainsi lorsque $f=ug$, où $u$ est inversible.

\textit{Proposition} \textbf{(1.1.10).} --- (i) \textit{Lorsque $f$ parcourt A, les ensembles $D(f)$ forment une base de la topologie de X.}

(ii) \textit{Pour tout $f\in A$, $D(f)$ est quasi compact.} \inserted{En particulier $X=D(1)$ est qu-cpt}

(i) Soit U un ensemble ouvert dans X; par définition, on a $U=X-V(E)$ où E est une partie de A, et $V(E)=\bigcap_{f\in E}V(f)$, d'où $U=\bigcup_{f\in E}D(f)$.

